\documentclass[12pt]{article}
\usepackage[T1]{fontenc}
\usepackage[utf8]{inputenc}
\usepackage{lmodern}
\usepackage{textcomp}
\usepackage{enumitem}
\usepackage{geometry}
\geometry{margin=1in}
\begin{document}
\begin{center}
Answer Key - Science - K-12 Level Assessment 3 \
\small (Answers are ordered exactly as in the question paper.)
\end{center}

\section*{Multiple Choice}
\begin{enumerate}[label=\arabic*.]
\item \textbf{Answer:} B
\\ \textit{Options:} A) high altitude; B) intense pressure; C) quick cooling; D) low temperatures
\\ \textit{Note:} Diamonds are formed from carbon under conditions of intense pressure and high temperature deep within the Earth's mantle, which facilitates the crystallization process necessary for diamond formation. This requirement for extreme pressure distinguishes diamond formation from other carbon-based materials.
\item \textbf{Answer:} A
\\ \textit{Options:} A) the same kind of cells working together.; B) different systems working together.; C) the same kind of organs forming a system.; D) different organs performing a function.
\\ \textit{Note:} The correct answer is appropriate because heart muscle, or cardiac muscle, is composed of specialized cells called cardiomyocytes that work collaboratively to contract and pump blood effectively throughout the body, highlighting the importance of cellular unity in its function.
\item \textbf{Answer:} A
\\ \textit{Options:} A) They can exist in the form of atoms or molecules.; B) They consist of atoms of two or more elements bonded together.; C) They have properties that are different from their component elements.; D) They can be broken down into elements by chemical means but not physical means.
\\ \textit{Note:} The correct answer is accurate because compounds are defined as substances formed from two or more different elements that are chemically bonded, specifically in molecule form, rather than existing as individual atoms. Therefore, stating that compounds can exist as atoms contradicts the fundamental definition of what a compound is.
\item \textbf{Answer:} C
\\ \textit{Options:} A) an increase in the size of the copper particles; B) a decrease in the mass of the copper particles; C) an increase in the motion of the copper particles; D) a decrease in the distance between the copper particles
\\ \textit{Note:} The correct answer is based on the principle that an increase in temperature causes particles, including those in a copper wire, to move more rapidly due to increased thermal energy, leading to heightened kinetic energy and motion among the copper particles. This heightened motion is what manifests as a temperature increase in the material.
\item \textbf{Answer:} B
\\ \textit{Options:} A) axial tilt; B) orbital path; C) the masses of the planets; D) the number of moons per planet
\\ \textit{Note:} The orbital path of the planets is most influenced by the gravitational force of the Sun because this force is what keeps the planets in their elliptical orbits, dictating their speed and trajectory as they revolve around the Sun. This gravitational interaction is fundamental to the structure and dynamics of the solar system.
\end{enumerate}

\section*{True / False}
\begin{enumerate}[label=\arabic*.]
\setcounter{enumi}{5}
\item \textbf{Answer:} False
\\ \textit{Options:} A) True; B) False
\\ \textit{Note:} Hot air warming a room is an example of heat transfer by convection, not radiation, as it involves the movement of air rather than electromagnetic waves.
\item \textbf{Answer:} False
\\ \textit{Options:} A) True; B) False
\\ \textit{Note:} Nerve cells, or neurons, do not divide to create new cells; instead, they are specialized cells that typically do not undergo mitosis, which is the process of cell division.
\item \textbf{Answer:} True
\\ \textit{Options:} A) True; B) False
\\ \textit{Note:} The statement is true because mechanical energy is the energy associated with the motion and position of an object, and when a nail is hammered, the kinetic energy from the hammer's movement is transferred to the nail, driving it into a surface. This physical movement exemplifies the transformation of mechanical energy in action.
\item \textbf{Answer:} False
\\ \textit{Options:} A) True; B) False
\\ \textit{Note:} Gemini is usually visible in the evening sky during January, not after midnight, as it sets earlier in the night.
\item \textbf{Answer:} True
\\ \textit{Options:} A) True; B) False
\\ \textit{Note:} The distance from the Sun to Earth is about 93 million miles, which is equivalent to approximately 500 light-seconds, making the statement true.
\end{enumerate}

\section*{Long Form Response}
\begin{enumerate}[label=\arabic*.]
\setcounter{enumi}{10}
\item \textbf{Answer:} Evidence that could suggest a bird species last seen in the 1900 s may still exist includes recent sightings, audio recordings of its calls, or the discovery of nests or feathers. This evidence challenges the claim of extinction by indicating that individuals of the species may still be alive in hidden or unexplored habitats. Researchers would need to investigate these findings further to confirm the bird's continued existence.
\\ \textit{Note:} correct because it identifies specific types of evidence-such as recent sightings, audio recordings, and physical signs like nests or feathers-that suggest the possibility of the bird species still existing, thereby challenging the notion of its extinction by implying that it could be surviving undetected in unexplored areas.
\item \textbf{Answer:} Preventing natural forest fires in a climax fire ecosystem can lead to an accumulation of dead biomass, which increases the risk of larger, more destructive fires in the future. This disruption can negatively impact the stability of the ecosystem by altering species composition and reducing biodiversity, as some plant and animal species rely on fire for regeneration and habitat maintenance. Over time, the ecosystem may become less resilient to pests and diseases, further threatening its health and diversity.
\\ \textit{Note:} correct because it highlights that preventing natural forest fires leads to an accumulation of dead biomass, which increases the likelihood of more severe fires, disrupts species composition, and diminishes biodiversity, ultimately reducing the ecosystem's resilience to threats like pests and diseases.
\end{enumerate}

\vfill
\noindent\textit{End of Answer Key}
\end{document}